% =====================================================================
%  SN4 (Supplementary Notes): AI-track review outcomes.
%  Referenced from the main text (sec:workflow-flags) as Supplementary Note SN4.
% =====================================================================


\section{AI-track review outcomes}
\label{sec:supp-ai-review}

This note reports the review outcomes for the two AI-drafted annotations:
Duerr et al.\ 2006, a conventional genome-wide association example, and
Inouye et al.\ 2018, a polygenic-score case that exposes several model
limitations.
% Review 2026-09-02: (1) 'interacting with the AI assistant under the review
% skill' restored: SN2 and SN3 describe the assistant as executing the skill
% under curator supervision, so 'the curator, who used the review skill'
% inverted the actor/tool relation; 'against' keeps the protocol phrase
% attached to 'reviewed', and the tie keeps 'v1.0' on the same line.
% (2) The human-authority sentence is restored verbatim from Section 5 and
% SN2, which the CHANGELOG records as one shared formulation. (3)
% 'authoritative logs' -> 'complete review logs' to avoid the
% authority/authoritative echo across the paragraph break.
Each draft was reviewed assertion by assertion against
\texttt{REVIEW\_PROTOCOL}~v1.0 by the human curator, who interacted with
an AI assistant operating under the review skill.
% Review 2026-09-03: the semicolon join is restored in the human-authority
% sentence: Section 5, SN2, and SN4 carry it as one verbatim shared
% formulation, recorded as such in the CHANGELOG and in the 2026-09-02 note
% above; the period split made SN4 silently diverge from the other two
% copies. The reordered first sentence of the paragraph is kept.
AI-generated drafts and
AI-assisted review suggestions were subject to final adjudication by the
human curator; the reviewed annotations are human-approved records, and no
AI system was treated as the authority for any scientific or curation
decision.

The complete review logs, including every flag and ruling, are available in
the repository under \fpath{annotations/reviews/}.
\Cref{tab:ai-review} summarizes the verdicts for the drafted items.

\begin{table}[H]  % [H] like every other supplement table
  \centering
  \small
  \begin{tabular}{lcccccc}
    \toprule
    Paper & Drafted & Approved & With flags & Edited & Rejected & Added on review \\
    \midrule
    Duerr 2006   & 6 & 0 & 2 & 4 & 0 & 1 (approved) \\
    Inouye 2018  & 3 & 1 & 1 & 1 & 0 & 1 (edited) \\
    \midrule
    Total        & 9 & 1 & 3 & 5 & 0 & 2 \\
    \bottomrule
  \end{tabular}
  \caption[Per-paper AI-track review verdicts]{Per-paper review verdicts
  for AI-drafted items. ``With flags'' denotes approval with non-blocking
  flags and no edits. ``Added on review'' denotes evidence omitted from the
  draft and entered by the reviewer as a new item.}
  \label{tab:ai-review}
\end{table}

The Duerr results reflect the \texttt{v0} annotation. The Inouye results,
and all related analyses in this note, use the \texttt{v1} re-annotation
produced under the protocol refined during the Duerr review. The original
five-item \texttt{v0} Inouye draft is retained in the repository for
comparison.

The drafts were produced under different protocol versions. The later Inouye
draft required fewer edits: two of its three items were accepted without
edits and one was edited, whereas four of Duerr's six items were edited.
This comparison is only suggestive, not controlled, because the papers differ
in both genre and difficulty. It is nevertheless consistent with the
provenance record: all thirty-seven cited Inouye spans were verbatim and
within the length limit, whereas the earlier Duerr draft contained notation
drift and one incorrect page citation. These issues are retained as
forward-looking notes because they predate the current rules.

A second pattern concerns omitted evidence. The reviewer-added items were not
granular association findings, which the drafts captured well. Instead, they
were higher-level or interpretive items: in Inouye, the paper's headline
claim that combining component scores outperforms each component; and in
Duerr, a cited anti-p40 antibody trial supporting the therapeutic-target
argument. The drafts were therefore most reliable for specific, directly
stated claims and less reliable for synthetic conclusions and cited
translational evidence that required cross-item interpretation.

Edits to existing items primarily concerned schema-level judgment rather than
the factual content of assertions. They included Knowledge Domain assignment,
distinguishing family-based from allele-frequency evidence; Method placement,
distinguishing targeted replication from a new genome-wide scan; Credibility
calibration, lowering the score-combination item from high to medium because
the paper gave no stated mechanism; and the applicability of Variant
Ascertainment to a polygenic-score target. In the last case, the reviewer
marked the dimension as not applicable and recorded the resulting constraint
as support for a candidate extension. No item was rejected, and no item was
found to misrepresent its source.

The revision also affects the promotion record. Because Variant Ascertainment
is no longer exercised by Inouye, its promotion rests on its initial proposal
in Davis and its independent reuse in Nelson and Duerr, which still satisfies
the two-papers rule. Updating the promotion record is a separate curation
step and is not undertaken here. Both reviews recommend a second pass over
the changed items because each paper gained a new item and several existing
items were edited.
